\documentclass[11pt]{article}
\usepackage[margin=1in]{geometry}
\usepackage{amsmath,amssymb,amsthm}
\usepackage{graphicx}
\usepackage[hidelinks]{hyperref}
\usepackage{microtype}

\newtheorem{theorem}{Theorem}
\newtheorem{proposition}{Proposition}
\newtheorem{corollary}{Corollary}
\newtheorem{lemma}{Lemma}

\newcommand{\R}{\mathbb R}
\newcommand{\Fl}{\operatorname{Fl}}
\newcommand{\depth}{\operatorname{depth}}

\title{A Gaussian--Section Obstruction to Uniform Common Depth\\
\large A negative range for Problem 2.3 of arXiv:2401.15095}
\author{Autonomous proof packet for \texttt{fa\_banach\_001}}
\date{17 August 2026}

\begin{document}
\maketitle

\noindent\textbf{Status.}
Candidate partial result, likely valid, pending human review.  The result gives
a broad negative parameter range, not a full classification of all tuples in
the source problem.

\section{Source problem}

Let $n=n_1+\cdots+n_r$ and let $\nu_i$ be the rank
$k=n_1+\cdots+n_i$ tautological bundle over the real flag manifold
$\Fl(n_1,\ldots,n_r)$.  Problem 2.3 of Antol\'in Camarena and Calles
Loperena asks for all tuples for which every collection of $m$ mass
assignments on $\nu_i$ admits a fiber and a point having ``sufficient depth''
for all $m$ measures \cite[Problem 2.3, PDF p.~4]{source}.  The preceding
paragraph explains that this means a depth bounded below by a positive
constant depending only on the base.

\begin{center}
\includegraphics[width=0.98\textwidth]{figures/open_problem_crop.png}
\end{center}

For a probability measure $\mu$ on a Euclidean space, vanishing on affine
hyperplanes, recall
\[
 \depth_{\mu}(x)=\inf\{\mu(H):H\text{ is a closed halfspace containing }x\}.
\]

\section{Idea of the proof}

A nowhere-zero section of $E^{\oplus(m-1)}$ is a continuous choice of
$m$ fiber points, with the first at the origin, that are not all equal.
After normalization, two of these centers stay a definite distance apart in
every fiber.  Put very narrow spherical Gaussians at the centers.  Every
candidate common point is then far from at least one center, and the Tukey
depth of the corresponding Gaussian is uniformly tiny.  Thus a topological
section becomes an adversarial family of perfectly regular mass assignments.

For the flag tautological bundle, $n-k+1$ fixed ambient vectors always give
such a section by orthogonal projection: they cannot all be perpendicular to
a $k$-plane.  A Stiefel--Whitney calculation shows that fewer summands cannot
have a nowhere-zero section, making the sectional threshold exact.

\section{The general bundle obstruction}

Let $E\to B$ be a rank-$k$ Euclidean vector bundle.  Define its best universal
simultaneous depth constant for $m$ assignments by
\[
 D_m(E)=\inf_{(\mu_1,\ldots,\mu_m)}
 \sup_{b\in B,\,x\in E_b}\min_{1\leq j\leq m}
 \depth_{\mu_j^b}(x),
\]
where the infimum runs over continuous mass assignments in the sense of the
source.

\begin{lemma}[spherical Gaussian depth]\label{lem:gauss}
Let $\gamma_{c,\sigma}$ be the Gaussian probability measure on $\R^k$ with
mean $c$ and covariance $\sigma^2I$.  If $\Phi$ is the standard normal
distribution function, then
\[
 \depth_{\gamma_{c,\sigma}}(x)
 =\Phi\!\left(-\frac{\lVert x-c\rVert}{\sigma}\right).
\]
\end{lemma}

\begin{proof}
For a unit vector $u$, the least-mass halfspace with inward normal $u$ that
contains $x$ has boundary through $x$.  Its mass on the side away from $c$ is
\[
 \Phi\!\left(-\frac{\langle x-c,u\rangle}{\sigma}\right).
\]
Minimizing over $u$ is the same as maximizing
$\langle x-c,u\rangle$, whose maximum is $\lVert x-c\rVert$.
\end{proof}

\begin{theorem}[Gaussian--section obstruction]\label{thm:general}
If $E^{\oplus(m-1)}$ has a nowhere-zero continuous section, then
\[
 D_m(E)=0.
\]
Equivalently, for every $\eta>0$ there are $m$ mass assignments such that no
fiber point has depth at least $\eta$ for all of them.
\end{theorem}

\begin{proof}
Write a nowhere-zero section as
\[
 S(b)=(s_2(b),\ldots,s_m(b))\in E_b^{\oplus(m-1)}
\]
and normalize it so that
\[
 \sum_{j=2}^m\lVert s_j(b)\rVert^2=1
 \qquad(b\in B).
\]
Set $c_1(b)=0$ and $c_j(b)=s_j(b)$ for $2\leq j\leq m$.  For each $b$, some
$j\geq2$ satisfies
$\lVert c_j(b)-c_1(b)\rVert\geq1/\sqrt{m-1}$.
Consequently, for every $x\in E_b$, the triangle inequality gives
\[
 \max_{1\leq a\leq m}\lVert x-c_a(b)\rVert
 \geq \frac{1}{2\sqrt{m-1}}.
\]

For $\sigma>0$, let $\mu_j^b$ be the spherical Gaussian on $E_b$ with center
$c_j(b)$ and covariance $\sigma^2I_{E_b}$.  In an orthonormal bundle chart,
the Gaussian integrals of bounded continuous test functions depend
continuously on the center and the fiber frame, so these are continuous mass
assignments.  They vanish on hyperplanes.

By Lemma~\ref{lem:gauss}, for every $b$ and $x$,
\[
 \min_j\depth_{\mu_j^b}(x)
 \leq
 \Phi\!\left(-\frac{1}{2\sigma\sqrt{m-1}}\right).
\]
The right side tends to zero as $\sigma\downarrow0$.  Given $\eta>0$, choose
$\sigma$ making it smaller than $\eta$.  This proves the claim.
\end{proof}

\section{Exact sectional threshold on the flag manifold}

Let $B=\Fl(n_1,\ldots,n_r)$, let $\nu_i$ have fiber $V_i$, and put
\[
 n=n_1+\cdots+n_r,
 \qquad k=n_1+\cdots+n_i<n.
\]

\begin{proposition}\label{prop:threshold}
The least positive integer $q$ for which $\nu_i^{\oplus q}$ admits a
nowhere-zero section is
\[
 q=n-k+1.
\]
\end{proposition}

\begin{proof}
Choose a fixed subspace $W\subset\R^n$ of dimension $n-k+1$, with orthonormal
basis $e_1,\ldots,e_{n-k+1}$.  At a flag with $k$-plane $V_i$, define
\[
 s_a(V_i)=P_{V_i}e_a\in V_i,
 \qquad 1\leq a\leq n-k+1.
\]
These sections cannot vanish simultaneously: otherwise
$W\subseteq V_i^\perp$, contradicting
$\dim W=n-k+1>n-k=\dim V_i^\perp$.  Hence their tuple is a nowhere-zero
section of $\nu_i^{\oplus(n-k+1)}$.

For the converse, the projection
\[
 \Fl(n_1,\ldots,n_r)\longrightarrow G_k(\R^n),
 \qquad (V_1,\ldots,V_{r-1})\longmapsto V_i,
\]
pulls the Grassmannian tautological bundle back to $\nu_i$ and is injective in
mod-$2$ cohomology, as in the projection used in the source's proof of its
main theorem \cite[Section 3]{source}.  The standard Grassmannian calculation
used there gives
\[
 w_k(\nu_i)^{n-k}\neq0.
\]
For $1\leq q\leq n-k$, this implies $w_k(\nu_i)^q\neq0$.  By the Whitney
product formula, the top Stiefel--Whitney class of the rank-$qk$ bundle
$\nu_i^{\oplus q}$ is
\[
 w_{qk}(\nu_i^{\oplus q})=w_k(\nu_i)^q\neq0.
\]
A bundle with a nowhere-zero section has zero top Stiefel--Whitney class.
Thus no smaller $q$ works.
\end{proof}

\begin{corollary}[negative range for Problem 2.3]\label{cor:flag}
If
\[
 \boxed{\qquad m\geq n-k+2,\qquad}
\]
then for every $\eta>0$ there are $m$ mass assignments on $\nu_i$ such that
no flag $\mathcal F$ and point $x\in V_i$ give depth at least $\eta$ for all
$m$ measures.  Hence no positive ``sufficient depth'' constant exists in this
parameter range.
\end{corollary}

\begin{proof}
Here $m-1\geq n-k+1$.  Pad the nowhere-zero section of
$\nu_i^{\oplus(n-k+1)}$ from Proposition~\ref{prop:threshold} with zero
components to obtain one of $\nu_i^{\oplus(m-1)}$, and apply
Theorem~\ref{thm:general}.
\end{proof}

\section{Verification, scope, and novelty}

The proof is exact and uses no numerical computation.  It gives a negative
classification region complementary to the source's positive
characteristic-class theorem.  The intermediate range
$2\leq m\leq n-k+1$ remains open in general; the absence of a nowhere-zero
center section there does not itself imply a positive common-depth theorem.

A bounded run-index and arXiv/web search on 17 August 2026 used arXiv id
2401.15095 and combinations of ``mass assignments'', ``Tukey depth'',
``Gaussian centers'', ``nowhere-zero section'', and ``$n-k+2$''.  It found
the source and related positive mass-partition results, but no exact negative
statement matching Theorem~\ref{thm:general} or Corollary~\ref{cor:flag}.
Novelty confidence is moderate.

\section*{Human-review recommendation}

Review is recommended.  The main checks are the exact Gaussian-depth formula,
weak continuity of the fiberwise Gaussians, and the mod-$2$ cohomology
injection/nonvanishing used to prove that the sectional threshold is exact.
The negative corollary itself only needs the explicit projection section, so
it remains valid even if the exact-minimality sentence requires adjustment.

\begin{thebibliography}{9}
\bibitem{source}
Omar Antol\'in Camarena and Jaime Calles Loperena,
\emph{A Center Transversal Theorem for mass assignments},
arXiv:2401.15095, 2024.
\end{thebibliography}

\end{document}
